\section{Discussion and limitations}
\label{sec:discussion}

\paragraph{What the benchmark says about implicit feedback.} Two numbers have to meet for ErrP feedback to help
a learner: the AUROC a decoder reaches for a new person, and the AUROC the learner needs. ErrP-Bench measures
the first across paradigms; the frontier measures the second. For observed robot errors, the paradigm that
motivates the field, they do not meet in our study: the best cross-subject decoder (\BestErrpHriRobot{}) is far
below the break-even level on reaching (\RLBreakEvenReach{}), and feedback at that quality slowed learning
relative to the sparse task reward. This agrees with a preliminary experiment of ours in which feedback from a
real cross-subject decoder of robot errors gave no gain over the sparse reward on the same task. It suggests that progress will come from paradigms and interfaces that make
errors more salient to the observer, or from per-user calibration, rather than from larger pretrained decoders.

\paragraph{Limitations.}
(i) The frontier uses synthetic feedback of controlled AUROC under an equal-variance binormal model with
independent noise per segment; real decoder errors are correlated across trials and participants, which could
move the break-even point. (ii) Two simulated tasks and one RL algorithm with a fixed feedback weight
($w=1$, frozen beforehand) are studied; a learner that discounts unreliable feedback could lower the
break-even AUROC. (iii) We did not tune hyperparameters for any decoder; tuning could favour the larger
models, although it would also break the equal-budget comparison. (iv) The foundation-model results depend on
our fine-tuning recipe (full fine-tuning, a new linear head, a short schedule). (v) EEGPT, the slowest model,
was run with fewer seeds and conditions than the others (Table~\ref{tab:main}). (vi) In IntEr-HRI, a motor
response falls inside most error windows; the early-window control shows that much, but not necessarily all,
of the signal precedes it. (vii) Cross-subject evaluation is our primary protocol; within-subject
calibration would reach higher AUROC and is not covered by the frontier placement.

\paragraph{Ethics and data.} All datasets are public recordings collected by their authors under their own
ethics approvals; we redistribute no EEG data, only preparation code. The ErrP-HRI dataset declares no
license; the others are released under Creative Commons licenses (BNCI 2015-013: CC BY-NC-ND 4.0; ERP CORE:
CC BY-SA 4.0; IntEr-HRI: CC BY 4.0).
